\chapter{视觉伺服}
\label{ch:visual-servoing}

% ════════════════════════════════════════════════════════════════
%  机械臂建模、规划与控制部 · 视觉伺服章
%  集大成融合：Siciliano《Modelling, Planning and Control》Ch10（系统全面：
%  交互矩阵/图像雅可比、PnP 位姿估计、PBVS/IBVS/2½-D 三族、对比表）
%  + Spong《Robot Modeling and Control》2e Ch11（工程直觉：交互矩阵逐步推导、
%  Lyapunov 控制律、camera-retreat 病态、分块法、运动可感知度）。
%  承上：相机模型（\cref{ch:camera}）给针孔投影/归一化坐标/内参，本章不重推；
%        机械臂雅可比/DLS（\cref{ch:arm-kin}）给 J 与阻尼伪逆，本章串联 (L_s J)；
%        操作空间控制（\cref{sec:actrl-osc}、\cref{ch:operational-space}）给笛卡尔控制底座。
%  记号归一：图像特征 s in R^k；交互矩阵 L_s（k×6，ṡ=L_s v_c）；
%        相机 twist v_c=[v;omega]（平移在前，与本书 ξ=[rho;φ] 同序）；增益 lambda；
%        几何雅可比 J（\cref{def:ak-jacobian}，注意其 twist=[omega;v] 角速度在前，串联需换序）。
%  归一化坐标取本书 camera_model 记号：s=[x,y]^T=[X/Z,Y/Z]，深度 Z。
% ════════════════════════════════════════════════════════════════

\begin{practice}[前置自测]
开始之前，先试答下列问题。能流畅作答者可速读本章表格与速查卡；卡住之处，正是本章要解决的。
\begin{enumerate}
  \item 一台相机装在机械臂末端，要它“看着”桌上一个零件并把末端移到抓取位姿。你会把误差定义在\emph{图像里}（像素坐标差）还是\emph{笛卡尔空间里}（位姿差）？两种选择各有什么代价？
  \item 相机以速度 $\mathbf{v}_c=[\mathbf{v};\boldsymbol{\omega}]$ 运动时，一个固定空间点在图像中的投影 $\mathbf{s}=[x,y]^\top$ 怎么动？这个 $\dot{\mathbf{s}}$ 与 $\mathbf{v}_c$ 的关系是线性的吗？系数矩阵依赖什么——只依赖图像坐标，还是还要知道点的深度 $Z$？
  \item 一个图像点给两个标量 $(x,y)$，相机有 $6$ 个自由度。那么“只盯着一个点不变”的相机运动有几维？这与上一章雅可比的\emph{零空间}是同一回事吗？
  \item 若要相机绕光轴旋转 $180^\circ$ 而图像里各特征点走直线，会发生什么灾难性的相机运动？为什么纯图像空间控制\emph{无法预测}相机轨迹？
  \item 把“期望的相机速度 $\mathbf{v}_c$”落到“关节速度 $\dot{\mathbf{q}}$”，需要把交互矩阵 $\mathbf{L}_s$ 与机械臂几何雅可比 $\mathbf{J}$（\cref{def:ak-jacobian}）怎样串起来？串联矩阵奇异时怎么办（回忆 \cref{thm:ak-dls}）？
  \item 相机装在手上（eye-in-hand）与装在场景里看着手（eye-to-hand），交互矩阵的符号会差一项吗？手眼之间那个未知刚体变换 $\mathbf{X}$ 怎么标出来（$\mathbf{A}\mathbf{X}=\mathbf{X}\mathbf{B}$）？
\end{enumerate}
\end{practice}

\section*{本章目标}
学完本章，你应当能够：
\begin{enumerate}
  \item \textbf{建立}视觉伺服的总框架：把相机测得的\emph{图像特征} $\mathbf{s}$ 反馈进控制，区分图像空间（IBVS）与笛卡尔空间（PBVS）两条主线及其混合（2½-D）；
  \item \textbf{推导}点特征的 $2\times6$ \emph{交互矩阵} $\mathbf{L}_s$（$\dot{\mathbf{s}}=\mathbf{L}_s\mathbf{v}_c$），理解深度 $Z$ 只进平移块、转动块只依赖图像坐标，并把 $n$ 个点堆叠成 $2n\times6$ 复合交互矩阵；
  \item \textbf{设计} IBVS 控制律 $\mathbf{v}_c=-\lambda\mathbf{L}_s^{+}(\mathbf{s}-\mathbf{s}^\ast)$，用 Lyapunov 分析其稳定性（$\mathbf{L}_s\widehat{\mathbf{L}}_s^{+}\succ0$），并解释深度估计、局部极小与\emph{相机后退}（camera retreat）病态；
  \item \textbf{设计} PBVS：由 PnP 估计位姿 $\mathbf{T}_o^c$，在操作空间做解析雅可比控制，并与 IBVS 在轨迹、鲁棒性、视野约束三方面对比；
  \item \textbf{掌握} 2½-D（混合）视觉伺服如何用单应分解把误差拆成“图像分量 + 笛卡尔分量”，兼得两者之长，并能填出三族方法的利弊对照表；
  \item \textbf{串联}视觉与机器人：把 $\mathbf{v}_c$ 经交互矩阵与几何雅可比 $\mathbf{J}$ 映到关节速度 $\dot{\mathbf{q}}=(\mathbf{L}_s\mathbf{J}_c)^{+}\dot{\mathbf{s}}$，奇异处用阻尼最小二乘（\cref{thm:ak-dls}），并区分 eye-in-hand 与 eye-to-hand；
  \item \textbf{求解}手眼标定 $\mathbf{A}\mathbf{X}=\mathbf{X}\mathbf{B}$（Tsai--Lenz）：闭式（旋转先于平移）与迭代解，并把它与相机内参标定（\cref{ch:calib}）衔接成完整的“从像素到关节力矩”的标定链。
\end{enumerate}

\subsection*{本章知识导航}
本章是一条“\emph{把相机看到的东西闭环回机器人运动}”的主线。它把两个已建好的支柱焊在一起：相机成像模型（\cref{ch:camera}：针孔投影、归一化坐标、内参）告诉我们\emph{世界如何变成像素}；机械臂雅可比（\cref{ch:arm-kin}：$\dot{\mathbf{x}}=\mathbf{J}\dot{\mathbf{q}}$、奇异、DLS）告诉我们\emph{关节如何驱动末端}。本章在二者之间补上缺失的中段——\emph{交互矩阵} $\mathbf{L}_s$，它把相机运动映到图像特征变化，从而让“图像误差”可以直接（IBVS）或间接（PBVS）地驱动关节。结构如下：

\begin{center}
\small
\begin{tikzpicture}[
  node distance=4mm and 8mm, >=Stealth,
  blk/.style={draw, rounded corners, align=center, font=\small, inner sep=3pt, fill=main!6, draw=main!60!black},
  grp/.style={draw, rounded corners, align=center, font=\small, inner sep=3pt, fill=second!6, draw=second!60!black},
  ln/.style={->, thick, gray!60!black}]
\node[blk] (feat) {图像特征 $\mathbf{s}$\\（点/线/矩）};
\node[blk, right=of feat] (L) {交互矩阵 $\mathbf{L}_s$\\ $\dot{\mathbf{s}}=\mathbf{L}_s\mathbf{v}_c$};
\node[grp, below=8mm of feat] (ibvs) {IBVS（图像空间）\\ $\mathbf{v}_c=-\lambda\mathbf{L}_s^{+}\mathbf{e}_s$};
\node[grp, right=of ibvs] (pbvs) {PBVS（位姿空间）\\ PnP $\to$ 笛卡尔控制};
\node[grp, right=of pbvs] (half) {2½-D（混合）\\ 单应分解};
\node[blk, below=8mm of ibvs, fill=third!8, draw=third!60!black] (robot) {接机器人 $(\mathbf{L}_s\mathbf{J}_c)^{+}$\\ 奇异$\to$DLS};
\node[blk, right=of robot, fill=third!8, draw=third!60!black] (handeye) {手眼标定\\ $\mathbf{A}\mathbf{X}=\mathbf{X}\mathbf{B}$};
\draw[ln] (feat)--(L);
\draw[ln] (L.south) -- ++(0,-4mm) -| (pbvs.north);
\draw[ln] (feat.south) -- (ibvs.north);
\draw[ln] (ibvs)--(pbvs); \draw[ln] (pbvs)--(half);
\draw[ln] (ibvs.south) -- (robot.north -| ibvs.south);
\draw[ln] (robot)--(handeye);
\end{tikzpicture}
\end{center}

\noindent\textbf{推荐路径}：第 \ref{sec:vs-overview} 节给总框架与两书的角度对照；第 \ref{sec:vs-features} 节回顾图像特征（承 \cref{ch:camera}，不重推相机模型）；第 \ref{sec:vs-interaction} 节是全章核心——交互矩阵的完整推导，必读；第 \ref{sec:vs-ibvs}、\ref{sec:vs-pbvs} 节分别讲 IBVS 与 PBVS；第 \ref{sec:vs-hybrid} 节讲 2½-D 并给三族对照表；第 \ref{sec:vs-robot} 节把 $\mathbf{v}_c$ 落到关节、讲 eye-in-hand vs eye-to-hand；第 \ref{sec:vs-handeye} 节讲手眼标定 $\mathbf{A}\mathbf{X}=\mathbf{X}\mathbf{B}$，与 \cref{ch:calib} 的内参标定合成完整标定链。

\paragraph{承上：本章站在两根柱子上。}
\cref{ch:camera} 给了针孔模型 $Z[u,v,1]^\top=\mathbf{K}(\mathbf{R}\mathbf{P}_w+\mathbf{t})$、归一化坐标 $[x,y,1]^\top=[X/Z,Y/Z,1]^\top$（小写 $x,y$ 在 $z=1$ 平面、无量纲）与内参 $\mathbf{K}$——本章把这些视为\emph{已知}，只用归一化坐标 $(x,y)$ 与深度 $Z$ 作图像特征，绝不重推成像几何。\cref{ch:arm-kin} 给了几何雅可比 $\mathbf{J}$（\cref{def:ak-jacobian}）、奇异与可操作度、阻尼最小二乘伪逆（\cref{thm:ak-dls}）——本章把相机速度 $\mathbf{v}_c$ 经 $\mathbf{J}$ 落到关节速度，并在奇异处复用 DLS。两本正典在此互补：Siciliano\cite{siciliano2009robotics} 系统全面（交互矩阵的一般形式、PnP、三族方法与严格 Lyapunov 证明），Spong\cite{spong2006robot} 工程直觉强（交互矩阵逐步代数推导、camera-retreat 病态的图像与几何、分块法）。本章\textbf{集二者之大成}而非二选一。

% ════════════════════════════════════════════════════════════════
\section{视觉伺服：总框架与两条主线}
\label{sec:vs-overview}

\paragraph{动机：把“看见”闭环成“动作”。}
视觉给机器人一种远距、非接触的环境感知：相机一帧图像就是一个二维光强矩阵，从中可实时提取\emph{图像特征参数}（feature parameters）——点的像素坐标、线的方向、区域的矩等。基于这些视觉量的反馈控制，称为\textbf{视觉伺服}（visual servoing）\cite{hutchinson1996tutorial,chaumette2006visual}。它与前几章的运动控制、力控的关键差别在于：\emph{被控量并非传感器直接测得}，而是要经图像处理与计算视觉从光强矩阵中提取、甚至进一步反演出来\cite{siciliano2009robotics}。

\begin{definition}[视觉伺服问题]\label{def:vs-problem}
设相机观测一个（暂设固定于基座系的）物体，从图像提取特征向量 $\mathbf{s}\in\mathbb{R}^k$。给定\emph{期望特征} $\mathbf{s}^\ast$（对应期望的相机--物体相对位姿），视觉伺服的目标是设计控制，使机器人末端（携带相机或被相机观测）运动，令特征误差
\begin{equation}
  \mathbf{e}_s=\mathbf{s}-\mathbf{s}^\ast
  \label{eq:vs-error}
\end{equation}
渐近趋零（$\mathbf{s}^\ast$ 可为常值——\emph{定点调节}，或时变——\emph{轨迹跟踪}）。
\end{definition}

\paragraph{两条主线：误差定义在哪里。}
机器人任务天然定义在\emph{操作空间}（末端相对物体的位姿），而视觉测量天然在\emph{图像平面}。把误差放在哪一侧，分出两族方法\cite{siciliano2009robotics,hutchinson1996tutorial}：

\begin{itemize}
  \item \textbf{位置/位姿空间视觉伺服}（position-based visual servoing, \textbf{PBVS}，亦称\emph{操作空间}视觉伺服）：先用视觉测量\emph{重建}物体相对相机的位姿 $\mathbf{T}_o^c$（解析 PnP 或迭代法），再把误差定义在笛卡尔空间，用操作空间控制驱动。优点：直接作用于位姿变量，笛卡尔轨迹可整定、可预测；缺点：依赖三维重建与精确标定，对标定误差敏感，且不直接控制图像，物体可能在过渡过程中\emph{逸出视野}导致反馈断路。
  \item \textbf{图像空间视觉伺服}（image-based visual servoing, \textbf{IBVS}）：直接以图像特征误差 \cref{eq:vs-error} 构造控制律，不显式重建三维位姿。优点：无需实时位姿估计，量在像素单位、对标定误差鲁棒（误差落在前向通道、抑制能力强），且直接控图像、易保持目标在视野内；缺点：图像与操作空间之间映射非线性，可能遇奇异、局部极小，末端笛卡尔轨迹\emph{无法事先预测}，可能撞障碍或越关节限位。
\end{itemize}
此外还有\textbf{混合}（2½-D）方法：误差一部分放图像、一部分放笛卡尔，取两者之长（\cref{sec:vs-hybrid}）。

\begin{insight}[一个矩阵把两条主线连起来]
无论 IBVS 还是 PBVS，核心都绕不开一个线性映射：相机速度 $\mathbf{v}_c$ 如何引起图像特征变化 $\dot{\mathbf{s}}$。这个映射的系数矩阵就是\textbf{交互矩阵} $\mathbf{L}_s$（\cref{sec:vs-interaction}）。IBVS 直接\emph{反用} $\mathbf{L}_s$（由想要的 $\dot{\mathbf{s}}$ 解出 $\mathbf{v}_c$）；PBVS 则用 $\mathbf{L}_s$（或其在多点上的版本）做\emph{数值位姿估计}（把“图像误差”反演成“位姿误差”，是 PnP 的迭代解法）。所以交互矩阵是全章的枢纽——正如几何雅可比 $\mathbf{J}$ 是 \cref{ch:arm-kin} 的枢纽。两者还高度类比：$\mathbf{J}$ 把关节速度映到末端 twist，$\mathbf{L}_s$ 把相机 twist 映到图像特征速度。
\end{insight}

\paragraph{相机放哪：eye-in-hand vs eye-to-hand。}
单相机系统有两种安装\cite{siciliano2009robotics}：\textbf{eye-in-hand}（手眼，相机固连末端，随臂运动）与 \textbf{eye-to-hand}（眼看手，相机固定于环境、观测末端与物体）。前者视野随末端走、近处分辨率高但易遮挡；后者全局视野稳定、但工作空间远端精度低。多相机系统还可混合二者。本章主体以 eye-in-hand、且\emph{末端系与相机系重合}的最简设定展开推导（这是两书的共同约定），eye-to-hand 的差异留到 \cref{sec:vs-robot} 专门处理（核心是交互矩阵到关节的映射差一项几何关系）。

% ════════════════════════════════════════════════════════════════
\section{图像特征回顾：点、线、矩}
\label{sec:vs-features}

\paragraph{动机：从几兆字节到几个数。}
一帧图像动辄数兆字节，但其\emph{信息量}远小于此。视觉伺服几乎总是先从图像提取少量\textbf{特征}：易检测、对尺度/朝向/光照有一定不变性、且在图像中可唯一识别\cite{spong2006robot}。好的特征应满足两点\cite{spong2006robot}：(i) 携带相机相对环境位姿的几何信息，相机运动与特征属性变化有良定关系（这是\cref{sec:vs-interaction} 交互矩阵存在的前提）；(ii) 易在图像序列中检测与跟踪。

\begin{note}[承 \cref{ch:camera}：成像几何不在此重推]\label{note:vs-feature-recall}
本章用到的成像关系全部来自 \cref{ch:camera}，此处只摘要、不重推：
\begin{itemize}
  \item \textbf{透视投影}：相机系点 $\mathbf{P}_c=[X,Y,Z]^\top$ 投影到归一化平面 $[x,y,1]^\top=[X/Z,\,Y/Z,\,1]^\top$（\cref{sec:cam-pinhole}）；再经内参 $\mathbf{K}$ 得像素 $[u,v,1]^\top=\mathbf{K}[x,y,1]^\top$（\cref{sec:cam-intrinsics}）。本章特征一律用\emph{归一化坐标} $\mathbf{s}=[x,y]^\top$（无量纲），因为它把内参剥离、使交互矩阵只含几何量；像素坐标只需在测量端用 $\mathbf{K}^{-1}$ 换算（要求内参已标定，\cref{ch:calib}）。
  \item \textbf{深度的角色}：透视除法丢失深度——单点像素只定一条\emph{视觉射线}（\cref{sec:cam-pinhole}），深度 $Z$ 需另行获取（双目 \cref{sec:cam-stereo}、RGB-D \cref{sec:cam-rgbd}、或估计）。这一点在交互矩阵里直接体现为“平移块含 $1/Z$”。
  \item \textbf{多视图几何}：两视图间的对极约束、本质/单应矩阵（\cref{sec:cam-epipolar}）是 PBVS 位姿估计与 2½-D 单应分解的工具，下文按需调用。
\end{itemize}
\end{note}

\paragraph{三类常用特征。}
\begin{itemize}
  \item \textbf{点特征}（最常用）：物体上一点在图像的投影坐标 $\mathbf{s}=[x,y]^\top$。角点（corner）定位性好——沿任意方向移动小窗口图像都明显变化，这是 Harris、Shi--Tomasi 等检测子的基础\cite{spong2006robot}；边缘（edge）只在垂直方向定位好、沿边方向滑动几乎不变，故不宜单独作伺服特征。$n$ 个点给 $2n$ 个特征分量，是本章推导的主角。
  \item \textbf{线特征}：图像中直线段，可用 $(\bar x,\bar y,L,\alpha)$（中点、长度、倾角）或 $(\rho,\theta)$ 参数化\cite{siciliano2009robotics}。其交互矩阵可由两端点的点交互矩阵导出。
  \item \textbf{矩特征}：在图像区域 $\mathcal{R}$ 上定义的矩（moments），可刻画二维目标的位置、朝向与形状\cite{siciliano2009robotics}。基于面积、质心、惯性比等的组合特征能\emph{解耦}相机自由度、改善 IBVS 的几何行为（如用面积控 $z$ 平移、用倾角控 $z$ 转动，见 \cref{sec:vs-ibvs} 的分块法）。
\end{itemize}

\begin{note}[特征检测与跟踪：本章的边界]\label{note:vs-tracking}
特征\emph{检测}（无先验地在图中找特征）与\emph{跟踪}（已知上帧位置、在本帧找同一特征）是两个相关但不同的问题\cite{spong2006robot}。给定（哪怕不确定的）相机运动模型，跟踪可用 Kalman/扩展 Kalman 滤波等估计方法处理。这些算法在开源视觉库中已有成熟实现，其细节超出本章范围——本章\emph{假定}特征已被可靠提取与配准，聚焦于“特征 $\to$ 相机运动”的控制理论。
\end{note}

% ════════════════════════════════════════════════════════════════
\section{交互矩阵（图像雅可比）}
\label{sec:vs-interaction}

\paragraph{动机：相机一动，图像怎么变。}
这是全章的数学核心。一旦选定特征 $\mathbf{s}$ 并能跟踪它，就要把\emph{相机运动}与\emph{特征变化率}关联起来。这关系是线性的（在给定构型下），其系数矩阵称\textbf{交互矩阵}（interaction matrix）或\textbf{图像雅可比}（image Jacobian）。它与机械臂雅可比（\cref{ch:arm-kin}）完全类比：都是把某个 $6$ 维速度（这里是相机 twist）线性映到一组参数的导数。

\subsection{定义与一般形式}
\label{sec:vs-interaction-def}

\begin{definition}[交互矩阵]\label{def:vs-interaction}
设特征向量 $\mathbf{s}\in\mathbb{R}^k$，相机相对固定物体的运动由\textbf{相机 twist}（在相机系下表达）
\begin{equation}
  \mathbf{v}_c=\begin{bmatrix}\mathbf{v}\\\boldsymbol{\omega}\end{bmatrix}\in\mathbb{R}^6,
  \qquad \mathbf{v}=\mathbf{R}_c^\top\dot{\mathbf{o}}_c,\quad \boldsymbol{\omega}=\mathbf{R}_c^\top\boldsymbol{\omega}_c
  \label{eq:vs-camera-twist}
\end{equation}
刻画（平移在前、转动在后，与本书 $\se(3)$ 排序 $\boldsymbol{\xi}=[\boldsymbol{\rho};\boldsymbol{\phi}]$ 同序，见 \cref{ch:lie}）。则存在 $k\times6$ 矩阵 $\mathbf{L}_s$，使
\begin{equation}
  \boxed{\dot{\mathbf{s}}=\mathbf{L}_s(\mathbf{s},Z)\,\mathbf{v}_c.}
  \label{eq:vs-interaction-def}
\end{equation}
$\mathbf{L}_s$ 称\textbf{交互矩阵}。它一般依赖当前特征值 $\mathbf{s}$ 与（点的）深度信息 $Z$。
\end{definition}

\begin{note}[交互矩阵 vs 图像雅可比：一个细微但重要的区分]\label{note:vs-Ls-vs-Js}
Siciliano\cite{siciliano2009robotics} 严格区分两者。\textbf{图像雅可比} $\mathbf{J}_s$ 把\emph{物体相对相机的相对 twist} $\mathbf{v}_{c,o}^c$ 映到 $\dot{\mathbf{s}}$；而\textbf{交互矩阵} $\mathbf{L}_s$ 把\emph{相机的绝对 twist} $\mathbf{v}_c^c$ 映到 $\dot{\mathbf{s}}$（设物体固定 $\mathbf{v}_o^c=\mathbf{0}$）。一般 $\dot{\mathbf{s}}=\mathbf{J}_s\mathbf{v}_o^c+\mathbf{L}_s\mathbf{v}_c^c$，两者经一个与相对位置有关的变换 $\boldsymbol{\Gamma}(\mathbf{o}_{c,o}^c)=\big[\begin{smallmatrix}-\mathbf{I}&\mathbf{S}(\cdot)\\\mathbf{0}&-\mathbf{I}\end{smallmatrix}\big]$ 相联：$\mathbf{L}_s=\mathbf{J}_s\boldsymbol{\Gamma}(\mathbf{o}_{c,o}^c)$。Spong\cite{spong2006robot} 则把二者统称“interaction matrix / image Jacobian”而不细分，因为伺服时常设物体固定、相机运动，此时只需 $\mathbf{L}_s$。\textbf{本章主用交互矩阵 $\mathbf{L}_s$}（物体固定、相机动），其解析式也更简洁；需要物体也运动（如传送带跟踪）时按上式补 $\mathbf{J}_s\mathbf{v}_o^c$ 项。
\end{note}

\subsection{点特征的 \texorpdfstring{$2\times6$}{2x6} 交互矩阵：完整推导}
\label{sec:vs-interaction-point}

下面对“运动相机观测固定空间点”这一最常见情形，完整推导其 $2\times6$ 交互矩阵。推导三步：(1) 求固定点相对运动相机系的速度；(2) 用透视投影把它转成图像速度；(3) 整理成矩阵。两书殊途同归，这里取 Siciliano 的紧凑微分法\cite{siciliano2009robotics}为主、并在符号与正负号上与 Spong 的逐步代数法\cite{spong2006robot}交叉核对。

\paragraph{第一步：固定点相对运动相机的速度。}
设空间点 $P$ 在基座系下位置 $\mathbf{p}$ 为常值，相机系原点 $\mathbf{o}_c$、姿态 $\mathbf{R}_c$。点在相机系下的坐标
\begin{equation}
  \mathbf{r}_c^c=\mathbf{R}_c^\top(\mathbf{p}-\mathbf{o}_c)=[x_c,\,y_c,\,z_c]^\top.
  \label{eq:vs-point-cam}
\end{equation}
对时间求导（$\mathbf{p}$ 常值），用 Poisson 公式 $\dot{\mathbf{R}}_c=\mathbf{R}_c\,\boldsymbol{\omega}^\wedge$（机体角速度，见 \cref{ch:rigid_body}）整理：
\begin{equation}
  \dot{\mathbf{r}}_c^c
  =-\mathbf{R}_c^\top\dot{\mathbf{o}}_c+\mathbf{S}(\mathbf{r}_c^c)\,\mathbf{R}_c^\top\boldsymbol{\omega}_c
  =\begin{bmatrix}-\mathbf{I}&\mathbf{S}(\mathbf{r}_c^c)\end{bmatrix}\mathbf{v}_c,
  \label{eq:vs-rdot}
\end{equation}
其中 $\mathbf{S}(\cdot)$ 是反对称（叉乘）矩阵。逐分量写出即 Spong 形式 $\dot{x}_c=z_c\omega_y-y_c\omega_z-\mathrm{v}_x$ 等（注意是\emph{固定}点相对\emph{动}系，故整体取负号——“固定点相对动系的速度，恰是动点相对固定系速度的 $-1$ 倍”\cite{spong2006robot}）。

\paragraph{第二步：透视投影的微分。}
取归一化坐标为特征 $\mathbf{s}=[x,y]^\top$，由 \cref{ch:camera}：
\begin{equation}
  \mathbf{s}=\frac{1}{z_c}\begin{bmatrix}x_c\\y_c\end{bmatrix}=\begin{bmatrix}x\\y\end{bmatrix}.
  \label{eq:vs-proj}
\end{equation}
对 \cref{eq:vs-proj} 用商法则求导，得投影的雅可比（$\partial\mathbf{s}/\partial\mathbf{r}_c^c$）：
\begin{equation}
  \dot{\mathbf{s}}=\frac{\partial\mathbf{s}}{\partial\mathbf{r}_c^c}\,\dot{\mathbf{r}}_c^c,
  \qquad
  \frac{\partial\mathbf{s}}{\partial\mathbf{r}_c^c}
  =\frac{1}{z_c}\begin{bmatrix}1&0&-x_c/z_c\\0&1&-y_c/z_c\end{bmatrix}
  =\frac{1}{z_c}\begin{bmatrix}1&0&-x\\0&1&-y\end{bmatrix}.
  \label{eq:vs-projjac}
\end{equation}

\paragraph{第三步：合并。}
把 \cref{eq:vs-rdot} 代入 \cref{eq:vs-projjac}，记深度 $Z\equiv z_c$，整理得点特征交互矩阵：
\begin{theorem}[点特征的交互矩阵]\label{thm:vs-point-interaction}
归一化坐标 $\mathbf{s}=[x,y]^\top$、深度 $Z$ 的点特征，其 $2\times6$ 交互矩阵为\cite{siciliano2009robotics,chaumette2006visual}
\begin{equation}
  \boxed{\;
  \mathbf{L}_s(\mathbf{s},Z)=
  \begin{bmatrix}
    -\dfrac{1}{Z} & 0 & \dfrac{x}{Z} & xy & -(1+x^2) & y\\[8pt]
    0 & -\dfrac{1}{Z} & \dfrac{y}{Z} & 1+y^2 & -xy & -x
  \end{bmatrix}\;}
  \label{eq:vs-Ls-point}
\end{equation}
满足 $\dot{\mathbf{s}}=\mathbf{L}_s\mathbf{v}_c$，其中 $\mathbf{v}_c=[\mathbf{v};\boldsymbol{\omega}]=[\mathrm{v}_x,\mathrm{v}_y,\mathrm{v}_z,\omega_x,\omega_y,\omega_z]^\top$。
\end{theorem}

\begin{note}[与 Spong 像素形式的等价、与符号核对]\label{note:vs-spong-form}
Spong\cite{spong2006robot} 用\emph{像素}坐标 $(u,v)$ 与焦距 $\lambda$，得
$\mathbf{L}_p=\big[\begin{smallmatrix}-\lambda/z&0&u/z&uv/\lambda&-(\lambda^2+u^2)/\lambda&v\\0&-\lambda/z&v/z&(\lambda^2+v^2)/\lambda&-uv/\lambda&-u\end{smallmatrix}\big]$。
令 $u=\lambda x,\,v=\lambda y$（即用归一化坐标 $x=u/\lambda$），两式逐项一致：例如 $-\lambda/z\to-(1/Z)\cdot\lambda$ 在按归一化坐标度量 $\dot{x}=\dot u/\lambda$ 后变 $-1/Z$；$uv/\lambda\to\lambda^2xy/\lambda=\lambda xy\to xy$；$-(\lambda^2+u^2)/\lambda\to-\lambda(1+x^2)\to-(1+x^2)$。故 \cref{eq:vs-Ls-point} 与 Spong 式仅差“以归一化坐标还是像素坐标度量特征”这一标度，本质同一矩阵。本书统一用归一化坐标的 \cref{eq:vs-Ls-point}（内参无关、最简洁）。\textbf{正负号已核对}：取相机沿光轴前进（$\mathrm{v}_z>0$），点在图像中向外扩张（$\dot x$ 与 $x$ 同号，由第 $3$ 列 $x/Z>0$ 给出），符合直觉。
\end{note}

\begin{insight}[读懂交互矩阵：深度只进平移块]\label{insight:vs-depth}
把 \cref{eq:vs-Ls-point} 按列分成平移块 $\mathbf{L}_v$（前 $3$ 列）与转动块 $\mathbf{L}_\omega$（后 $3$ 列）：
\[
  \dot{\mathbf{s}}=\underbrace{\begin{bmatrix}-1/Z&0&x/Z\\0&-1/Z&y/Z\end{bmatrix}}_{\mathbf{L}_v(\mathbf{s},Z)}\mathbf{v}
  +\underbrace{\begin{bmatrix}xy&-(1+x^2)&y\\1+y^2&-xy&-x\end{bmatrix}}_{\mathbf{L}_\omega(\mathbf{s})}\boldsymbol{\omega}.
\]
\textbf{关键观察}\cite{spong2006robot,chaumette2006visual}：$\mathbf{L}_v$ 依赖深度 $Z$，$\mathbf{L}_\omega$ \emph{只依赖图像坐标}、与 $Z$ 无关。物理上很合理：纯转动相机不改变深度感（远近点的角位移规律相同），纯平移则有“近大远小”的视差——深度未知时，转动部分仍精确，平移部分只差一个标度。这正是 IBVS 对深度估计误差相对鲁棒的根源（\cref{sec:vs-ibvs}），也是分块法（partitioned method）的设计依据。
\end{insight}

\subsection{多点堆叠与零空间}
\label{sec:vs-interaction-multi}

\paragraph{多点：竖向堆叠。}
$n$ 个点的特征向量 $\mathbf{s}=[x_1,y_1,\dots,x_n,y_n]^\top\in\mathbb{R}^{2n}$，各点深度 $Z_i$。复合交互矩阵是各点 $\mathbf{L}_{s_i}$ 的竖向堆叠\cite{siciliano2009robotics,spong2006robot}：
\begin{equation}
  \dot{\mathbf{s}}=\mathbf{L}_s\mathbf{v}_c,\qquad
  \mathbf{L}_s=\begin{bmatrix}\mathbf{L}_{s_1}(\mathbf{s}_1,Z_1)\\\vdots\\\mathbf{L}_{s_n}(\mathbf{s}_n,Z_n)\end{bmatrix}\in\mathbb{R}^{2n\times6}.
  \label{eq:vs-Ls-stack}
\end{equation}
每点贡献 $2$ 行；要唯一确定 $6$ 维相机运动，至少需 $3$ 个点（$2n\ge6$）使 $\mathbf{L}_s$ 满秩。实践常用 $4$ 个（含冗余、抗噪、且共面 $4$ 点的位姿估计唯一，见 \cref{sec:vs-pbvs}）。

\begin{insight}[单点交互矩阵的零空间：四维不可观运动]\label{insight:vs-nullspace}
单点 $\mathbf{L}_s\in\mathbb{R}^{2\times6}$ 秩为 $2$，零空间四维：存在 $4$ 维相机运动不改变该点投影\cite{spong2006robot,siciliano2009robotics}。其中两维有直观意义：(i) 沿\emph{视觉射线}（过焦点与 $P$ 的直线）平移——点沿射线移近移远，投影不变；(ii) 绕该视觉射线\emph{转动}——投影不变。这与 \cref{ch:arm-kin} 的雅可比零空间是同一类概念：$\mathbf{L}_s$ 的零空间就是“相机能动但图像不变”的方向，正如 $\mathbf{J}$ 的零空间是“关节能动但末端不动”的方向。多点堆叠后零空间收缩——$\ge3$ 点时一般退化为零，相机任何运动都引起可观图像变化。
\end{insight}

% ════════════════════════════════════════════════════════════════
\section{图像空间视觉伺服（IBVS）}
\label{sec:vs-ibvs}

\paragraph{动机：直接把图像误差变成相机速度。}
IBVS 不重建三维位姿，而是把特征误差 $\mathbf{e}_s=\mathbf{s}-\mathbf{s}^\ast$ 直接映成相机运动。思路：想让误差按 $\dot{\mathbf{e}}_s=-\lambda\mathbf{e}_s$ 指数收敛，再用交互矩阵 $\dot{\mathbf{s}}=\mathbf{L}_s\mathbf{v}_c$ 反解出实现它的 $\mathbf{v}_c$。

\subsection{控制律与稳定性}
\label{sec:vs-ibvs-law}

\begin{theorem}[IBVS 比例控制律]\label{thm:vs-ibvs-law}
设物体固定、$\mathbf{s}^\ast$ 常值，故 $\dot{\mathbf{e}}_s=\dot{\mathbf{s}}=\mathbf{L}_s\mathbf{v}_c$。欲令 $\dot{\mathbf{e}}_s=-\lambda\mathbf{e}_s$（$\lambda>0$），由 $-\lambda\mathbf{e}_s=\mathbf{L}_s\mathbf{v}_c$ 解出相机速度命令
\begin{equation}
  \boxed{\;\mathbf{v}_c=-\lambda\,\mathbf{L}_s^{+}\,\mathbf{e}_s=-\lambda\,\mathbf{L}_s^{+}(\mathbf{s}-\mathbf{s}^\ast),\;}
  \label{eq:vs-ibvs-law}
\end{equation}
其中 $\mathbf{L}_s^{+}$ 为交互矩阵的 Moore--Penrose 伪逆：$k>6$（特征多于自由度，常见）取左伪逆 $\mathbf{L}_s^{+}=(\mathbf{L}_s^\top\mathbf{L}_s)^{-1}\mathbf{L}_s^\top$；$k<6$ 取右伪逆并可加零空间项 $\mathbf{v}_c=\mathbf{L}_s^{+}(-\lambda\mathbf{e}_s)+(\mathbf{I}-\mathbf{L}_s^{+}\mathbf{L}_s)\mathbf{b}$；$k=6$ 且满秩则 $\mathbf{L}_s^{+}=\mathbf{L}_s^{-1}$。
\end{theorem}

\begin{proof}[稳定性（Lyapunov）]
取 $V=\tfrac12\|\mathbf{e}_s\|^2\ge0$\cite{spong2006robot,hutchinson1996tutorial}。沿闭环 $\dot V=\mathbf{e}_s^\top\dot{\mathbf{e}}_s=\mathbf{e}_s^\top\mathbf{L}_s\mathbf{v}_c=-\lambda\,\mathbf{e}_s^\top\mathbf{L}_s\mathbf{L}_s^{+}\mathbf{e}_s$。
\begin{itemize}
  \item $k=6$ 满秩：$\mathbf{L}_s\mathbf{L}_s^{+}=\mathbf{I}$，$\dot V=-\lambda\|\mathbf{e}_s\|^2=-2\lambda V<0$，\emph{指数稳定}。
  \item $k>6$：$\mathbf{L}_s\mathbf{L}_s^{+}\in\mathbb{R}^{k\times k}$ 只半正定（秩 $\le6<k$，有非平凡零空间），故只能证\emph{稳定}、不能证渐近稳定——可能停在 $\mathbf{e}_s\neq\mathbf{0}$ 但 $\mathbf{e}_s\in\ker\mathbf{L}_s^\top$ 的\textbf{局部极小}（见下）。
\end{itemize}
\end{proof}

\begin{note}[深度未知与对标定误差的鲁棒性]\label{note:vs-depth-est}
$\mathbf{L}_s$ 含深度 $Z$，而 IBVS 哲学本不愿估计 $Z$。实践用\emph{估计}的 $\widehat{\mathbf{L}}_s$：可取期望构型的深度 $Z^\ast$、初始构型深度、或在线估计，控制律变 $\mathbf{v}_c=-\lambda\widehat{\mathbf{L}}_s^{+}\mathbf{e}_s$\cite{chaumette2006visual,siciliano2009robotics}。Lyapunov 分析给出鲁棒性判据：闭环稳定当 $\mathbf{L}_s\widehat{\mathbf{L}}_s^{+}\succ0$（正定）\cite{spong2006robot}。由于 $Z$ 的误差只标度平移块 $\mathbf{L}_v$（\cref{insight:vs-depth}），适度的深度误差不破坏正定性——这解释了 IBVS 对深度与内参标定误差的著名鲁棒性。常见三种 $\widehat{\mathbf{L}}_s$ 选法：当前 $\mathbf{L}_s(\mathbf{s},\widehat Z)$、期望 $\mathbf{L}_s(\mathbf{s}^\ast,Z^\ast)$、或两者均值 $\tfrac12(\mathbf{L}_s+\mathbf{L}_s^\ast)$，各有收敛域与轨迹形状的折中\cite{chaumette2006visual}。
\end{note}

\subsection{病态：相机后退与局部极小}
\label{sec:vs-ibvs-retreat}

\paragraph{相机后退（camera retreat）。}
IBVS 显式控制\emph{图像}误差，却对\emph{相机轨迹}毫无显式约束——这是其阿喀琉斯之踵\cite{spong2006robot,chaumette2006visual}。经典反例：图像平面与四个共面特征点所在平面平行，期望位姿与初始位姿仅差\emph{绕光轴的纯转动}。IBVS 令各特征点在图像里走\emph{直线}到目标，图像误差指数收敛；但为实现这种直线轨迹，相机被迫沿光轴\emph{后退}（图像点直线运动 + 绕光轴转，几何上只能靠后退兼顾）。

\begin{insight}[最极端情形：退到无穷远]\label{insight:vs-retreat-pi}
当所需相机运动是绕光轴\emph{转} $\pi$ 时，灾难达到极致\cite{spong2006robot}：各特征点直线轨迹\emph{都过图像中心}，这只在相机后退到 $z=-\infty$ 时才能发生——此时 $\det\mathbf{L}_s=0$，控制失效。根因在 \cref{insight:vs-depth}：转动块 $\mathbf{L}_\omega$ 与平移块 $\mathbf{L}_v$ 在“绕光轴转”与“沿光轴退”之间产生病态耦合。这是 IBVS 局部行为的标志性失败，也是 2½-D 与分块法要解决的首要问题。
\end{insight}

\begin{figure}[H]
\centering
\begin{tikzpicture}[scale=1.0,>=Stealth]
% IBVS: straight lines through center -> retreat
\begin{scope}[shift={(0,0)}]
  \draw[thick,gray!50] (-1.2,-1.2) rectangle (1.2,1.2);
  \fill[main!70!black] (1.0,0.6) circle (1.6pt); \fill[main!70!black] (-1.0,-0.6) circle (1.6pt);
  \fill[main!70!black] (-0.6,1.0) circle (1.6pt); \fill[main!70!black] (0.6,-1.0) circle (1.6pt);
  \draw[->,second!70!black,thick] (1.0,0.6)--(-1.0,-0.6);
  \draw[->,second!70!black,thick] (-0.6,1.0)--(0.6,-1.0);
  \node at (0,-1.65) {\scriptsize IBVS：直线过中心 $\Rightarrow$ 相机退至 $z=-\infty$};
\end{scope}
% partitioned/2.5D: circular arcs, no retreat
\begin{scope}[shift={(4.2,0)}]
  \draw[thick,gray!50] (-1.2,-1.2) rectangle (1.2,1.2);
  \fill[main!70!black] (1.0,0.6) circle (1.6pt); \fill[main!70!black] (-1.0,-0.6) circle (1.6pt);
  \draw[->,third!60!black,thick] (1.0,0.6) arc (31:211:1.166);
  \node at (0,-1.65) {\scriptsize 分块/2½-D：圆弧 $\Rightarrow$ 相机原地转、不后退};
\end{scope}
\end{tikzpicture}
\caption{相机后退病态与其化解。左：纯 IBVS 下各特征点走直线，绕光轴转 $\pi$ 时轨迹过图像中心，迫使相机退至 $z=-\infty$（$\det\mathbf{L}_s=0$）。右：分块法/2½-D 让特征点走圆弧，相机原地转动、不后退。}
\label{fig:vs-retreat}
\end{figure}

\paragraph{分块法（partitioned method）：一种补救。}
把 \cref{eq:vs-Ls-point} 按自由度拆开，用图像伺服控 $x,y$ 轴的平动与转动（$\mathbf{v}_{xy},\boldsymbol{\omega}_{xy}$），另用\emph{专门特征}控光轴自由度 $\mathrm{v}_z,\omega_z$\cite{spong2006robot}：以两特征点连线与图像横轴的夹角 $\theta_{ij}$ 控 $\omega_z=\gamma_{\omega_z}(\theta_{ij}^\ast-\theta_{ij})$；以图像中某多边形面积 $\sigma^2$ 控 $\mathrm{v}_z=\gamma_{\mathrm{v}_z}\ln(\sigma^\ast/\sigma)$（面积对绕光轴转动不变，故\emph{解耦}了转动与 $z$ 平移）。其余自由度由
\begin{equation}
  \boldsymbol{\xi}_{xy}=-\mathbf{L}_{xy}^{+}\big(\lambda\mathbf{e}_s+\mathbf{L}_z\boldsymbol{\xi}_z\big)
  \label{eq:vs-partitioned}
\end{equation}
给出，其中 $\mathbf{L}_z\boldsymbol{\xi}_z$ 是 $z$ 自由度引起的特征运动、被当作\emph{已知扰动}从误差里扣除。如此相机不再后退（\cref{fig:vs-retreat} 右），代价是图像误差不再单调下降。这是用“矩/几何特征解耦自由度”改良 IBVS 的范例，与用图像矩做特征（\cref{sec:vs-features}）一脉相承。

\begin{insight}[运动可感知度：图像版的可操作度]\label{insight:vs-perceptibility}
\cref{ch:arm-kin} 用可操作度 $w=\sqrt{\det(\mathbf{J}\mathbf{J}^\top)}$ 量化“关节速度 $\to$ 末端速度”的放大。视觉伺服有完全平行的概念——\textbf{运动可感知度}（motion perceptibility）$w_v=\sqrt{\det(\mathbf{L}_s^\top\mathbf{L}_s)}=\sigma_1\cdots\sigma_6$，量化“相机速度 $\to$ 图像特征速度”的放大\cite{spong2006robot}。单位相机速度球经 $\mathbf{L}_s$ 映成图像特征速度椭球，$w_v$ 正比其体积；$w_v\to0$ 处某方向相机运动几乎不引起图像变化（伺服在该方向“失明”、增益需提高、噪声放大）。把 \cref{ch:arm-kin} 的可操作度椭球直接搬到图像空间即得——又一处雅可比与交互矩阵的深层类比。
\end{insight}

% ════════════════════════════════════════════════════════════════
\section{位姿空间视觉伺服（PBVS）}
\label{sec:vs-pbvs}

\paragraph{动机：先看清位姿，再做笛卡尔控制。}
PBVS 把视觉伺服拆成两段：先用视觉\emph{估计}物体相对相机的位姿 $\mathbf{T}_o^c$（位姿估计/PnP），再把误差定义在笛卡尔空间、用操作空间控制（\cref{sec:actrl-osc}、\cref{ch:operational-space}）驱动。它在概念上等同于“以视觉为反馈源的操作空间控制”。

\subsection{位姿估计：PnP 与迭代解}
\label{sec:vs-pbvs-pnp}

\begin{definition}[PnP 问题]\label{def:vs-pnp}
\textbf{Perspective-$n$-Point}（PnP）：已知物体上 $n$ 个点在\emph{物体系}下的坐标 $\mathbf{r}_{o,i}^o$（如来自 CAD 模型）与它们在\emph{图像平面}的投影 $\mathbf{s}_i$，求物体相对相机的位姿 $\mathbf{T}_o^c$\cite{siciliano2009robotics}。每个对应给出 $\lambda_i\widetilde{\mathbf{s}}_i=\boldsymbol{\Pi}\mathbf{T}_o^c\widetilde{\mathbf{r}}_{o,i}^o$ 形式的约束。
\end{definition}

\begin{note}[解的个数：几何决定多义性]\label{note:vs-pnp-solutions}
PnP 解数依点数与共面性\cite{siciliano2009robotics}：
\begin{itemize}
  \item P3P（$3$ 非共线点）：\emph{四}解；
  \item P4P/P5P（非共面）：至少两解；但 $\ge4$ 共面点且无三点共线时\emph{唯一}；
  \item PnP（$n\ge6$ 非共面）：\emph{唯一}解。
\end{itemize}
这解释了 IBVS/PBVS 为何常取 $4$ 个共面点——既够估计、解又唯一。共面情形的解析解尤为简洁：令物体平面为 $z_o=0$，约束退化为可解的单应问题（用 \cref{sec:cam-epipolar} 的单应工具），由 $\mathbf{H}$ 分解出 $\mathbf{R}_o^c$ 与 $\mathbf{o}_{c,o}^c$。
\end{note}

\paragraph{迭代（数值）位姿估计：把 PnP 化成“逆运动学”。}
解析 PnP 较繁，常用\emph{迭代}法\cite{siciliano2009robotics}：把位姿估计当作一个虚拟伺服——用最小坐标 $\mathbf{x}_{c,o}$ 表示位姿，定义“当前估计位姿渲染出的虚拟特征”$\mathbf{s}(\widehat{\mathbf{x}}_{c,o})$，令其逼近实测 $\mathbf{s}$。误差 $\mathbf{s}-\mathbf{s}(\widehat{\mathbf{x}}_{c,o})$ 经解析图像雅可比 $\mathbf{J}_{A_s}$ 反演更新 $\widehat{\mathbf{x}}_{c,o}$：
\begin{equation}
  \dot{\widehat{\mathbf{x}}}_{c,o}=\mathbf{J}_{A_s}^{-1}(\widehat{\mathbf{x}}_{c,o})\,\mathbf{K}_s\big(\mathbf{s}-\mathbf{s}(\widehat{\mathbf{x}}_{c,o})\big),
  \label{eq:vs-virtual-vs}
\end{equation}
数值积分至收敛。此法（亦称\emph{虚拟视觉伺服}）与机械臂逆运动学的数值解完全同构——把“关节角 $\to$ 末端位姿”换成“位姿 $\to$ 图像特征”\cite{siciliano2009robotics}。其奇异既含姿态表示奇异、又含交互矩阵奇异（后者更关键，依赖特征选取）。用作\emph{视觉跟踪}时，上一帧估计作本帧初值，收敛快、精度高。

\subsection{PBVS 控制律}
\label{sec:vs-pbvs-law}

得到 $\mathbf{T}_o^c$ 与期望 $\mathbf{T}_o^d$ 后，构造当前相机系相对期望相机系的位姿偏差
\begin{equation}
  \mathbf{T}_c^d=\mathbf{T}_o^d(\mathbf{T}_o^c)^{-1}=\begin{bmatrix}\mathbf{R}_c^d&\mathbf{o}_{d,c}^d\\\mathbf{0}^\top&1\end{bmatrix},
  \label{eq:vs-pose-err}
\end{equation}
定义操作空间误差 $\widetilde{\mathbf{x}}=-[\,\mathbf{o}_{d,c}^d;\,\boldsymbol{\phi}_{d,c}\,]$（$\boldsymbol{\phi}_{d,c}$ 为 $\mathbf{R}_c^d$ 的欧拉角）\cite{siciliano2009robotics}。控制目标令 $\widetilde{\mathbf{x}}\to\mathbf{0}$。两种实现：

\begin{itemize}
  \item \textbf{PD + 重力补偿}（动力学层）：$\mathbf{u}=\mathbf{g}(\mathbf{q})+\mathbf{J}_{A_d}^\top(\mathbf{K}_P\widetilde{\mathbf{x}}-\mathbf{K}_D\mathbf{J}_{A_d}\dot{\mathbf{q}})$，与运动控制的 PD+重力补偿（\cref{sec:actrl-osc}）同构，仅误差定义不同；Lyapunov 函数 $V=\tfrac12\dot{\mathbf{q}}^\top\mathbf{B}\dot{\mathbf{q}}+\tfrac12\widetilde{\mathbf{x}}^\top\mathbf{K}_P\widetilde{\mathbf{x}}$ 证渐近稳定\cite{siciliano2009robotics}。
  \item \textbf{解析速度}（resolved-velocity，运动学层）：设臂有高增益内环、近似理想定位器 $\mathbf{q}\approx\mathbf{q}_r$，直接算关节参考速度 $\dot{\mathbf{q}}_r=\mathbf{J}_{A_d}^{-1}\mathbf{K}\widetilde{\mathbf{x}}$，得线性闭环 $\dot{\widetilde{\mathbf{x}}}+\mathbf{K}\widetilde{\mathbf{x}}=\mathbf{0}$，指数收敛\cite{siciliano2009robotics}。此法避开了视觉低帧率与动力学耦合，工程常用。$\mathbf{K}$ 取对角等增益时，相机原点走直线段——\emph{笛卡尔轨迹可预测}，正是 PBVS 相对 IBVS 的优势。
\end{itemize}

\begin{insight}[IBVS vs PBVS：误差放哪，代价就在哪]\label{insight:vs-ibvs-vs-pbvs}
两者是一组深刻的对偶\cite{siciliano2009robotics,chaumette2006visual}：
\begin{itemize}
  \item \textbf{轨迹}：PBVS 笛卡尔轨迹直、可预测（相机走直线），但图像轨迹不受控、目标可能逸出视野；IBVS 图像轨迹直、目标留在视野，但笛卡尔轨迹不可预测、可能后退或撞限位。
  \item \textbf{鲁棒性}：IBVS 量在像素、标定误差落\emph{前向}通道（抑制强）、且期望特征由相机实测，故对内/外参误差鲁棒；PBVS 标定误差经位姿重建落\emph{反馈}通道（抑制弱），更敏感。
  \item \textbf{物体几何}：单相机 PBVS 需已知物体几何（做 PnP）；IBVS 单相机原则上不需（仅 $Z$ 这一项与几何有关）。
  \item \textbf{失败模式}：IBVS 怕局部极小与相机后退；PBVS 怕目标逸出视野与位姿估计不收敛。
\end{itemize}
2½-D 正是想同时拿到“PBVS 的笛卡尔可预测”与“IBVS 的视野保持”。
\end{insight}

% ════════════════════════════════════════════════════════════════
\section{2½-D（混合）视觉伺服}
\label{sec:vs-hybrid}

\paragraph{动机：各取一半，规避两端的坑。}
2½-D 视觉伺服（Malis--Chaumette）把误差\emph{一半放笛卡尔、一半放图像}\cite{malis1999twohalf,siciliano2009robotics}：转动用从单应分解得的姿态误差（笛卡尔，保证相机姿态轨迹可预测），平移用图像特征 + 一个深度比（保证目标留在视野）。名字“2½-D”意指它介于纯 3D（PBVS）与纯 2D（IBVS）之间——既不需完整三维重建，又不纯靠图像。

\paragraph{构造。}
设物体平面上 $\ge4$ 点。由当前与期望图像的对应，算\emph{平面单应} $\mathbf{H}$（不需知当前/期望相机位姿，只需 $\mathbf{s},\mathbf{s}^\ast$）\cite{siciliano2009robotics}：
\begin{equation}
  \mathbf{H}=\mathbf{R}_d^c+\frac{1}{d_d}\mathbf{o}_{c,d}^c\,\mathbf{n}^{d\top},
  \label{eq:vs-homography}
\end{equation}
其中 $\mathbf{R}_d^c$ 为期望到当前的姿态、$\mathbf{n}^d$ 为平面法向、$d_d$ 为期望位姿到平面的距离。分解 $\mathbf{H}$ 得 $\mathbf{R}_d^c$、$\mathbf{n}^d$ 与 $(1/d_d)\mathbf{o}_{c,d}^c$。于是：
\begin{itemize}
  \item \textbf{转动控制}（笛卡尔）：由 $\mathbf{R}_d^c$ 取欧拉角 $\boldsymbol{\phi}_{c,d}$，令 $\boldsymbol{\omega}_r^c=-\mathbf{T}(\boldsymbol{\phi}_{c,d})\mathbf{K}_o\boldsymbol{\phi}_{c,d}$，得 $\dot{\boldsymbol{\phi}}_{c,d}+\mathbf{K}_o\boldsymbol{\phi}_{c,d}=\mathbf{0}$，指数收敛；
  \item \textbf{平移控制}（图像 + 深度比）：用某点的图像坐标差与深度比 $\rho_z=z_c/z_d$ 的对数构造位置误差
  \begin{equation}
    \mathbf{e}_p=\begin{bmatrix}x_d-x\\y_d-y\\\ln\rho_z\end{bmatrix},
    \label{eq:vs-hybrid-ep}
  \end{equation}
  其中 $\rho_z$ 可由单应的行列式等量算得（\cref{eq:vs-homography} 的副产物）。$\mathbf{e}_p\to\mathbf{0}$ 等价于 $\mathbf{r}_d^c-\mathbf{r}_c^c\to\mathbf{0}$。
\end{itemize}
如此，某选定点在图像中走\emph{完美直线}（由 $x,y$ 分量直接控），相机姿态轨迹又可预测，笛卡尔轨迹明显优于纯 IBVS（\cref{fig:vs-retreat} 右的圆弧即此类行为）\cite{siciliano2009robotics}。

\begin{table}[H]
\centering
\small
\caption{三族视觉伺服方法对照（集 Siciliano\cite{siciliano2009robotics} 与 Chaumette\cite{chaumette2006visual,chaumette2007visual} 之要点）}
\label{tab:vs-comparison}
\begin{tabular}{@{}llll@{}}
\toprule
 & \textbf{PBVS（位姿）} & \textbf{IBVS（图像）} & \textbf{2½-D（混合）}\\
\midrule
误差空间 & 笛卡尔位姿 & 图像特征 & 转动笛卡尔 + 平移图像\\
需位姿估计 & 是（PnP/迭代） & 否 & 部分（单应分解）\\
需物体几何 & 是（单目） & 否（除 $Z$） & 否（共面假设）\\
笛卡尔轨迹 & 直、可预测 & 不可预测、可能后退 & 较可预测\\
图像/视野 & 不受控、易逸出 & 直接控、易保持 & 选定点直线、较好\\
标定鲁棒性 & 较低（反馈通道） & 高（前向通道） & 中\\
奇异/极小 & 位姿估计不收敛 & 局部极小、camera retreat & 大幅缓解二者\\
代表文献 & \cite{siciliano2009robotics} & \cite{espiau1992new,hutchinson1996tutorial} & \cite{malis1999twohalf}\\
\bottomrule
\end{tabular}
\end{table}

\begin{note}[谱系：交互矩阵与三族方法的源流]\label{note:vs-lineage}
交互矩阵的概念由 Weiss 等早期工作提出、由 Espiau--Chaumette--Rives\cite{espiau1992new} 正式命名并系统化（“task function”框架）。Hutchinson--Hager--Corke 的经典综述\cite{hutchinson1996tutorial} 统一了 IBVS/PBVS 的术语。Malis--Chaumette--Boudet\cite{malis1999twohalf} 提出 2½-D。Chaumette--Hutchinson 的两部 tutorial\cite{chaumette2006visual,chaumette2007visual} 是现代视觉伺服的标准入门（I 讲基本方法、II 讲高级与几何特征）。本章把这条谱系的核心——交互矩阵、三族控制律、稳定性——融两书之长重述到复现级。
\end{note}

% ════════════════════════════════════════════════════════════════
\section{接机器人：从相机速度到关节速度}
\label{sec:vs-robot}

\paragraph{动机：$\mathbf{v}_c$ 不是终点，关节力矩才是。}
前几节的输出都是\emph{相机速度} $\mathbf{v}_c$。要真正驱动机械臂，必须把它落到关节速度 $\dot{\mathbf{q}}$，这一步把视觉（交互矩阵 $\mathbf{L}_s$）与机器人（几何雅可比 $\mathbf{J}$，\cref{def:ak-jacobian}）串联起来。

\subsection{eye-in-hand：串联交互矩阵与几何雅可比}
\label{sec:vs-robot-eih}

eye-in-hand 且相机系与末端系重合时，相机 twist 即末端 twist，由几何雅可比给出。但要小心\textbf{排序}：本书几何雅可比的 twist 是 $\boldsymbol{\mathcal{V}}=[\boldsymbol{\omega}_e;\mathbf{v}_e]$（\emph{角速度在前}，\cref{def:ak-jacobian}），而相机 twist $\mathbf{v}_c=[\mathbf{v};\boldsymbol{\omega}]$（\emph{平移在前}，\cref{eq:vs-camera-twist}）。二者差一个换序置换 $\mathbf{P}_{\text{swap}}=\big[\begin{smallmatrix}\mathbf{0}&\mathbf{I}_3\\\mathbf{I}_3&\mathbf{0}\end{smallmatrix}\big]$。定义“与相机同序”的雅可比 $\mathbf{J}_c=\mathbf{P}_{\text{swap}}\mathbf{J}$，则 $\mathbf{v}_c=\mathbf{J}_c\dot{\mathbf{q}}$。代入 $\dot{\mathbf{s}}=\mathbf{L}_s\mathbf{v}_c$：

\begin{theorem}[视觉--关节复合映射]\label{thm:vs-LsJ}
eye-in-hand、相机系=末端系时，图像特征速度与关节速度的复合映射为
\begin{equation}
  \dot{\mathbf{s}}=\mathbf{L}_s\,\mathbf{J}_c\,\dot{\mathbf{q}}=\mathbf{J}_L\,\dot{\mathbf{q}},
  \qquad \mathbf{J}_L\equiv\mathbf{L}_s\mathbf{J}_c\in\mathbb{R}^{k\times n}.
  \label{eq:vs-JL}
\end{equation}
反解得关节参考速度（resolved-velocity IBVS）
\begin{equation}
  \boxed{\;\dot{\mathbf{q}}=\mathbf{J}_L^{+}\,\mathbf{K}_s\,\mathbf{e}_s=(\mathbf{L}_s\mathbf{J}_c)^{+}\mathbf{K}_s(\mathbf{s}^\ast-\mathbf{s}).\;}
  \label{eq:vs-resolved-joint}
\end{equation}
工程上分两步算更稳：先由图像求相机速度 $\mathbf{v}_r^c=\mathbf{L}_s^{+}\mathbf{K}_s\mathbf{e}_s$，再由雅可比求关节速度 $\dot{\mathbf{q}}=\mathbf{J}_c^{-1}\mathbf{v}_r^c$\cite{siciliano2009robotics}。
\end{theorem}

\begin{insight}[双重奇异：交互矩阵奇异 + 机器人奇异]\label{insight:vs-double-singular}
复合映射 $\mathbf{J}_L=\mathbf{L}_s\mathbf{J}_c$ 在\emph{两类}构型退化\cite{siciliano2009robotics}：(i) 机械臂\textbf{运动学奇异}（$\mathbf{J}$ 掉秩，\cref{sec:ak-singularity}）——末端某方向不可动；(ii) 交互矩阵\textbf{奇异}（$\mathbf{L}_s$ 掉秩）——某相机运动图像不可观，依赖特征选取，更隐蔽也更关键。两类奇异处直接用伪逆都会让 $\dot{\mathbf{q}}$ 爆炸。补救与 \cref{ch:arm-kin} 完全一致：用\textbf{阻尼最小二乘}（DLS，\cref{thm:ak-dls}）
\begin{equation}
  \dot{\mathbf{q}}=\mathbf{J}_L^\top(\mathbf{J}_L\mathbf{J}_L^\top+\mu^2\mathbf{I})^{-1}\mathbf{K}_s\mathbf{e}_s
  \quad\text{或对 }\mathbf{L}_s\text{、}\mathbf{J}_c\text{ 分别加阻尼},
  \label{eq:vs-dls}
\end{equation}
以少量精度换数值稳定（奇异值 $\sigma\to\sigma/(\sigma^2+\mu^2)$，\cref{eq:ak-dls}）。对交互矩阵奇异，另一招是\emph{增加特征数} $k>6$ 用左伪逆，远离 $\mathbf{L}_s$ 的零空间\cite{siciliano2009robotics}。冗余臂还可用零空间投影 $\mathbf{N}=\mathbf{I}-\mathbf{J}_L^{+}\mathbf{J}_L$ 做次要任务（避障、避关节限位），与 \cref{sec:ak-singularity} 的冗余框架同构。
\end{insight}

\subsection{eye-to-hand：差一个固定变换}
\label{sec:vs-robot-eth}

eye-to-hand 时相机固定、观测随末端运动的物体（或末端本身）。此时\emph{物体}相对相机的运动由末端运动诱导，交互矩阵映射的速度是物体 twist，需经末端到物体的几何关系换算\cite{siciliano2009robotics,chaumette2007visual}。形式上：物体随末端刚体运动，其在相机系下的 twist 由末端 twist 经伴随变换 $\mathrm{Ad}$（\cref{ch:lie}）给出；与 eye-in-hand 相比，关键差别是\textbf{符号}与\textbf{参考系}——eye-in-hand 是“相机动、点固定”，eye-to-hand 是“相机固定、点随手动”，二者诱导的 $\dot{\mathbf{s}}$ 差一个整体符号（回忆 \cref{eq:vs-rdot} 的负号来自“固定点相对动系”）。实现上只需把 \cref{eq:vs-JL} 的 $\mathbf{J}_c$ 换成“末端运动 $\to$ 物体在相机系 twist”的正确雅可比（含固定的相机--基座外参与手--物体变换），其余控制律不变。

% ════════════════════════════════════════════════════════════════
\section{手眼标定 \texorpdfstring{$\mathbf{A}\mathbf{X}=\mathbf{X}\mathbf{B}$}{AX=XB}}
\label{sec:vs-handeye}

\paragraph{动机：相机装在手上，但“装在哪”得标出来。}
要把视觉测得的位姿用到机器人控制，必须知道\emph{相机系与末端系（或基座系）之间的固定刚体变换} $\mathbf{X}$——这是 eye-in-hand 系统的外参。\cref{ch:calib} 标的是相机\emph{内参}（焦距、主点、畸变），本节标的是这个\emph{手眼外参}，二者合起来才打通“像素 $\to$ 关节”的完整链条。经典提法是 Tsai--Lenz 的 $\mathbf{A}\mathbf{X}=\mathbf{X}\mathbf{B}$\cite{tsai1989hand}。

\subsection{问题构造}
\label{sec:vs-handeye-formulation}

让机器人末端带相机运动到若干位姿，每个位姿下相机观测同一固定标定物（其位姿由 \cref{ch:calib} 的标定板估计得到）。设两次运动间：
\begin{itemize}
  \item $\mathbf{A}$＝两次\emph{末端}位姿之间的相对变换（由机器人正运动学 \cref{ch:arm-kin} 精确读出，$\mathbf{A}=\mathbf{T}_{e_2}^{-1}\mathbf{T}_{e_1}$ 一类）；
  \item $\mathbf{B}$＝两次\emph{相机}相对标定物的位姿之间的相对变换（由视觉位姿估计得出）；
  \item $\mathbf{X}=\mathbf{T}_c^e$＝待求的手眼变换（末端系 $\to$ 相机系，固定不变）。
\end{itemize}
由“末端动 = 相机动经 $\mathbf{X}$ 共轭”得约束
\begin{equation}
  \boxed{\;\mathbf{A}\mathbf{X}=\mathbf{X}\mathbf{B},\;}
  \qquad \mathbf{A},\mathbf{B},\mathbf{X}\in\SE(3).
  \label{eq:vs-axxb}
\end{equation}
$N$ 次运动给 $N$ 个这样的方程，联立求唯一 $\mathbf{X}$。

\begin{note}[可解性条件]\label{note:vs-axxb-solvability}
单个 \cref{eq:vs-axxb} 不足以定 $\mathbf{X}$：至少需\emph{两次}转轴方向\emph{不平行}的相对运动\cite{tsai1989hand}。直觉上，绕单一轴的运动无法约束 $\mathbf{X}$ 沿该轴的分量。实践取 $\ge3$ 次姿态变化充分的运动以抗噪。
\end{note}

\subsection{闭式解：旋转先于平移（Tsai--Lenz）}
\label{sec:vs-handeye-closed}

把 \cref{eq:vs-axxb} 按旋转/平移分块（$\mathbf{A}=(\mathbf{R}_A,\mathbf{t}_A)$ 等），得\emph{解耦}的两组方程\cite{tsai1989hand}：
\begin{align}
  \mathbf{R}_A\mathbf{R}_X&=\mathbf{R}_X\mathbf{R}_B, \label{eq:vs-axxb-rot}\\
  (\mathbf{R}_A-\mathbf{I})\mathbf{t}_X&=\mathbf{R}_X\mathbf{t}_B-\mathbf{t}_A. \label{eq:vs-axxb-trans}
\end{align}
\textbf{先解旋转}：\cref{eq:vs-axxb-rot} 表明 $\mathbf{R}_A$ 与 $\mathbf{R}_B$ 是\emph{相似}的旋转，其转轴经 $\mathbf{R}_X$ 对应。用旋转的轴角/四元数表示：设 $\mathbf{R}_A,\mathbf{R}_B$ 的旋转轴（用 Rodrigues 向量或“修正 Rodrigues” $\mathbf{P}_A,\mathbf{P}_B$）满足
\begin{equation}
  \mathbf{R}_X\mathbf{P}_B=\mathbf{P}_A
  \;\Longrightarrow\;
  (\mathbf{P}_A+\mathbf{P}_B)^\wedge\,\mathbf{p}_X'=\mathbf{P}_A-\mathbf{P}_B,
  \label{eq:vs-tsai-rot}
\end{equation}
这是关于 $\mathbf{R}_X$ 参数 $\mathbf{p}_X'$ 的\emph{线性}方程，$\ge2$ 次运动堆叠后最小二乘解出，再由 $\mathbf{p}_X'$ 重构 $\mathbf{R}_X$\cite{tsai1989hand}。\textbf{再解平移}：$\mathbf{R}_X$ 已知后 \cref{eq:vs-axxb-trans} 对 $\mathbf{t}_X$ 是\emph{线性}的，堆叠各次运动的 $(\mathbf{R}_A-\mathbf{I})$ 与右端，最小二乘解出 $\mathbf{t}_X$。这种“旋转先、平移后”的解耦是 Tsai--Lenz 法的精髓——把一个 $\SE(3)$ 上的非线性问题拆成两个线性最小二乘。

\begin{note}[闭式 vs 迭代；以及 $\mathbf{A}\mathbf{X}=\mathbf{Z}\mathbf{B}$ 推广]\label{note:vs-handeye-variants}
\textbf{闭式}（分离旋转平移）快、无需初值，但旋转误差会传播进平移、且未联合优化\cite{tsai1989hand}。\textbf{迭代/同时}法在 $\SE(3)$ 上联合最小化所有方程残差（如用李代数 \cref{ch:lie} 参数化做非线性最小二乘，或四元数/对偶四元数同时解旋转平移），精度更高、对噪声更稳，代价是需初值（常以闭式解作初值）。\textbf{eye-to-hand} 的标定结构类似，但求的是相机--基座变换，方程形式推广为 $\mathbf{A}\mathbf{X}=\mathbf{Z}\mathbf{B}$（$\mathbf{X}$＝手--物体、$\mathbf{Z}$＝相机--基座，两个未知变换联立）。
\end{note}

\begin{insight}[标定链：内参 + 手眼 = 像素到关节]\label{insight:vs-calib-chain}
视觉伺服要落地，需要一条完整标定链\cite{tsai1989hand,siciliano2009robotics}：(1) 相机\textbf{内参}标定（\cref{ch:calib}：$\mathbf{K}$、畸变）——把像素换成归一化坐标，是本章一切交互矩阵的前提；(2) \textbf{手眼}外参标定（本节：$\mathbf{X}=\mathbf{T}_c^e$）——把相机系测得的位姿/速度换到末端系，是 \cref{sec:vs-robot} 串联 $\mathbf{L}_s$ 与 $\mathbf{J}$ 的前提。二者共享同一套数学（PnP/标定板位姿估计），且都用 \cref{ch:calib} 的标定板数据——手眼标定的 $\mathbf{B}$ 正来自内参标定时每帧估计的板位姿。这条链把 \cref{ch:camera}（成像）、\cref{ch:calib}（内参）、本章（伺服 + 手眼）焊成一个从“光子打到传感器”到“关节发出力矩”的闭环。
\end{insight}

% ════════════════════════════════════════════════════════════════
\section{本章小结}
\label{sec:vs-summary}

本章把\emph{视觉}闭环回\emph{机器人运动}，集 Siciliano\cite{siciliano2009robotics} 与 Spong\cite{spong2006robot} 两本正典之大成：

\begin{enumerate}
  \item \textbf{交互矩阵是枢纽}：点特征的 $2\times6$ 交互矩阵 \cref{eq:vs-Ls-point}（$\dot{\mathbf{s}}=\mathbf{L}_s\mathbf{v}_c$）把相机 twist 映到图像特征速度，与几何雅可比 $\mathbf{J}$ 完全类比。深度 $Z$ 只进平移块、转动块只依赖图像坐标（\cref{insight:vs-depth}）——这是 IBVS 对深度误差鲁棒的根源。
  \item \textbf{三族方法各有取舍}（\cref{tab:vs-comparison}）：IBVS 直接反用 $\mathbf{L}_s$（$\mathbf{v}_c=-\lambda\mathbf{L}_s^{+}\mathbf{e}_s$），对标定鲁棒、保视野，但有局部极小与相机后退病态（\cref{insight:vs-retreat-pi}）；PBVS 先 PnP 估位姿再做笛卡尔控制，轨迹可预测，但对标定敏感、目标易逸出视野；2½-D 用单应分解各取一半，缓解两端的坑。
  \item \textbf{接机器人}：相机速度经 $(\mathbf{L}_s\mathbf{J}_c)^{+}$ 落到关节速度（\cref{eq:vs-resolved-joint}），注意 twist 排序换序；双重奇异（交互矩阵 + 机器人）处用 DLS（\cref{thm:ak-dls}）稳健化；eye-in-hand 与 eye-to-hand 差一个固定变换与符号。
  \item \textbf{标定链闭环}：手眼标定 $\mathbf{A}\mathbf{X}=\mathbf{X}\mathbf{B}$（Tsai--Lenz，旋转先于平移的解耦闭式解）求出相机--末端外参，与 \cref{ch:calib} 的内参标定合成“像素到关节”的完整链条（\cref{insight:vs-calib-chain}）。
\end{enumerate}

\begin{note}[与全书的接口]\label{note:vs-interfaces}
本章上承 \cref{ch:camera}（成像几何、归一化坐标、内参）与 \cref{ch:arm-kin}（几何雅可比、奇异、DLS、冗余零空间），下接 \cref{ch:operational-space}（操作空间动力学/阻抗——PBVS 的笛卡尔控制底座，亦可把视觉伺服与力控融成“视觉--力”混合控制）。交互矩阵与雅可比、运动可感知度与可操作度的层层类比（\cref{insight:vs-perceptibility,insight:vs-double-singular}），是把本章嵌进全书“雅可比即枢纽”主线的纽带。更进一步的话题——基于图像矩的解耦特征、深度学习驱动的特征提取与端到端伺服、视觉--惯性--力多模态融合——可在此理论骨架上展开。
\end{note}

\begin{practice}[复盘]
\begin{enumerate}
  \item 由 \cref{eq:vs-rdot} 与 \cref{eq:vs-projjac} 亲手补全 \cref{eq:vs-Ls-point} 的推导，逐项核对符号（提示：先写出 $\dot x=\partial(x_c/z_c)/\partial t$）。
  \item 验证 \cref{note:vs-spong-form}：把 Spong 像素形式 $\mathbf{L}_p$ 用 $u=\lambda x,v=\lambda y$ 代入并按 $\dot x=\dot u/\lambda$ 度量，逐项还原到归一化形式 \cref{eq:vs-Ls-point}。
  \item 对单点 $\mathbf{L}_s$（\cref{eq:vs-Ls-point}）验证：向量 $[x,y,1,0,0,0]^\top$（沿视觉射线平移）与 $[0,0,0,x,y,1]^\top$（绕射线转）都在其零空间，解释 \cref{insight:vs-nullspace}。
  \item 解释为何相机绕光轴转 $\pi$ 会让纯 IBVS 退到 $z=-\infty$；分块法（\cref{eq:vs-partitioned}）如何用面积特征 $\sigma$ 与角度特征 $\theta_{ij}$ 化解（\cref{fig:vs-retreat}）。
  \item 对一个 6R 臂 + eye-in-hand 相机，写出从 $4$ 个图像点误差到关节速度的完整 resolved-velocity IBVS 流水线（含 $\mathbf{P}_{\text{swap}}$ 换序、$\widehat{\mathbf{L}}_s$ 深度估计、奇异处 DLS）。
  \item 设计一组手眼标定运动：为什么至少需两次\emph{转轴不平行}的运动（\cref{note:vs-axxb-solvability}）？给出 $\mathbf{A},\mathbf{B}$ 各自如何由机器人正运动学与视觉位姿估计获得。
\end{enumerate}
\end{practice}
